\chapter{Phase compactness and the perimeter lower bound}\label{sec:phase-compactness}
The energy bound alone forces two phases and a finite-perimeter interface; no Euler equation is needed. We give the details of \cite{ChAC}, Proposition 3.5 and Exercise 3.1, using the BV compactness already established in Theorem \ref{thm:bv-local-compactness}. Throughout this section $U\subset\mathbb R^n$ is bounded and open.

\begin{lemma}[Quartic coercivity and concentration at the wells]\label{lem:phase-coercivity}
For every real $s$,
\[
 W(s)\geq\frac18|s|^4-\frac14,
 \qquad \operatorname{dist}(s,\{-1,1\})^2\leq4W(s).
\]
Thus if $E_\varepsilon(u;U)\leq C$, then
\[
 \int_U|u|^4\leq8C\varepsilon+2|U|,
 \qquad
 \int_U\operatorname{dist}(u,\{-1,1\})^2\leq4C\varepsilon.
\]
In particular the measure where the distance from the wells exceeds $\delta>0$ is at most $4C\varepsilon/\delta^2$.
\end{lemma}
\begin{proof}
Subtracting $|s|^4/8-1/4$ from $W(s)$ gives $(s^2-2)^2/8$. Also $4W(s)=(|s|-1)^2(|s|+1)^2\geq(|s|-1)^2$. Integrate and use $\int W(u)\leq C\varepsilon$; Chebyshev's inequality gives the last statement. These are the coercivity and well-concentration estimates used in \cite{ChAC}, pp.~8--9.
\end{proof}

\begin{proposition}[Binary phases and perimeter normalization]\label{prop:phase-bv-normalization}
For a measurable set $E\subset U$, write $u_E=2\chi_E-1$. Then $u_E\in BV(U)$ if and only if $E$ has finite perimeter in $U$, and
\[
 Du_E=2D\chi_E,\qquad |Du_E|=2|D\chi_E|,
 \qquad D\Phi(u_E)=c_WD\chi_E.
\]
For two measurable sets, $\|u_E-u_F\|_{L^1(U)}=2|E\triangle F|$.
\end{proposition}
\begin{proof}
Differentiate the affine expressions in distributions; the constant function has zero derivative relative to $U$. Since $\Phi(u_E)=c_W\chi_E-c_W/2$, the last derivative identity follows. The distance identity is pointwise. This is the two-phase specialization of Example \ref{ex:bv-jump-coarea}; in particular $|D\Phi(u_E)|=(c_W/2)|Du_E|$, retaining the factor two in \cite{ChAC}, Proposition 3.5.
\end{proof}

\begin{theorem}[Compactness of bounded diffuse-interface energy]\label{thm:phase-compactness}
Let $\varepsilon_j\downarrow0$ and $u_j\in H^1(U)\cap L^4(U)$ satisfy $\sup_j E_{\varepsilon_j}(u_j;U)<\infty$. A subsequence converges in $L^p(U)$ for every $1\leq p<4$ to $u_E=2\chi_E-1$, where $E$ has finite perimeter in $U$. Along this subsequence $\Phi(u_j)\to\Phi(u_E)$ in $L^1_{\mathrm{loc}}(U)$ and
\[
 c_W P(E;U)\leq\liminf_j E_{\varepsilon_j}(u_j;U).
\]
No boundary regularity or boundary trace condition is assumed.
\end{theorem}
\begin{proof}
Lemma \ref{lem:phase-coercivity} bounds $u_j$ in $L^4(U)$. The cubic growth of $\Phi$ bounds $\Phi(u_j)$ in $L^1(U)$, and Lemma \ref{lem:phase-modica-energy} bounds their total variations. BV compactness and a diagonal almost-everywhere extraction give $\Phi(u_j)\to v$ in $L^1_{\mathrm{loc}}$ and almost everywhere. Continuity of $\Phi^{-1}$ gives $u_j\to u_0=\Phi^{-1}(v)$ almost everywhere. Fatou and $\int W(u_j)\to0$ show $u_0\in\{-1,1\}$ almost everywhere.

The $L^4$ bound gives uniform integrability on the finite-measure set $U$. More explicitly, $\int_A|u_j|\leq C|A|^{3/4}$, and truncation at a large value makes the remaining tails uniformly small. Dominated convergence for the bounded truncations then gives $L^1(U)$ convergence. Interpolation between this convergence and the uniform $L^4$ bound gives convergence in every $L^p$, $p<4$.

Write $u_0=u_E$. Lower semicontinuity of variation applied to $\Phi(u_j)$ gives $c_WP(E;U)\leq\liminf_j\int_U|D\Phi(u_j)|\leq\liminf_jE_{\varepsilon_j}(u_j;U)$. Here the lower-semicontinuity theorem permits local $L^1$ convergence even when the variation is measured on all of $U$. The finite volume gives $u_E\in L^1(U)$, so $E$ has finite relative perimeter. This fills out \cite{ChAC}, Proposition 3.5; quartic coercivity upgrades the local convergence there to global convergence on bounded $U$.
\end{proof}

\begin{theorem}[The perimeter lower bound for a prescribed limit]\label{thm:phase-liminf}
Suppose $\varepsilon_j\downarrow0$ and $u_j\to u$ in $L^1(U)$. Then
\[
 \liminf_j E_{\varepsilon_j}(u_j;U)\geq
 \begin{cases}
 c_WP(E;U),&u=2\chi_E-1\text{ a.e.},\ E\text{ of finite perimeter},\\
 +\infty,&\text{otherwise}.
 \end{cases}
\]
If the energies are uniformly bounded and $u=u_E$, the localized bound
$c_WP(E;A)\leq\liminf_j E_{\varepsilon_j}(u_j;A)$ holds for every open $A\subset U$.
\end{theorem}
\begin{proof}
If the liminf is finite, select a subsequence attaining it with bounded energy and apply Theorem \ref{thm:phase-compactness}. Its phase limit must be the prescribed $L^1$ limit. This proves the global assertion, including infinite perimeter. For any fixed open $A$, restrict the sequence to $A$ and repeat the same argument for its energy liminf; $A$ is bounded and the prescribed convergence restricts to it. The subsequence used to prove the inequality may depend on $A$, but the inequality concerns the original sequence. This is the liminf conclusion of \cite{ChAC}, Proposition 3.5.
\end{proof}

\begin{proposition}[The limiting energy measure dominates perimeter]\label{prop:phase-measure-bound}
Assume the hypotheses of Theorem \ref{thm:phase-compactness}, with $u_j\to u_E$ in $L^1(U)$. If a subsequence of the energy measures converges weakly on compactly supported continuous tests in $U$ to $\mu$, then
\[
 \mu\geq c_W|D\chi_E|
\]
as Radon measures on $U$.
\end{proposition}
\begin{proof}
Pass further to the primitive-convergent subsequence in Theorem \ref{thm:phase-compactness}; this preserves $\mu$ and $u_E$. For $X\in C_c^1(U;\mathbb R^n)$, integration by parts and the pointwise energy inequality give
\[
 \left|\int_U\Phi(u_j)\operatorname{div}X\right|
 \leq\int_U|X|\,d\mu_{\varepsilon_j}^{u_j}.
\]
Pass to the limit. The resulting bound says that the vector measure $D\Phi(u_E)$ has total variation at most $\mu$, by the dual characterization of total variation with localized test fields. Proposition \ref{prop:phase-bv-normalization} identifies this variation as $c_W|D\chi_E|$. Equality is an additional assertion and need not hold.
\end{proof}

\begin{corollary}[Energy convergence forces equipartition]\label{cor:phase-equipartition}
Suppose $u_j\to u_E$ in $L^1(U)$, $\varepsilon_j\downarrow0$, and
$E_{\varepsilon_j}(u_j;U)\to c_WP(E;U)<\infty$. Then
\[
 \int_U|D\Phi(u_j)|\to c_WP(E;U),\qquad
 |\xi_{\varepsilon_j}^{u_j}|(U)\to0.
\]
Each of the gradient and potential energies tends to $c_WP(E;U)/2$.
\end{corollary}
\begin{proof}
For every subsequence, bounded energy gives a further subsequence whose primitives converge locally to $\Phi(u_E)$. Lower semicontinuity gives the lower bound $c_WP(E;U)$ for their variations, while the energy inequality gives the matching upper bound. The subsequence criterion therefore gives convergence along the whole sequence. Its energy defect tends to zero. Proposition \ref{prop:phase-defect} makes the total variation of discrepancy tend to zero. Since the sum of the two energies converges to $c_WP(E;U)$ and their difference tends to zero, each converges to half the sum.
\end{proof}

\begin{example}[Two disappearing layers can retain energy]\label{ex:phase-hidden-double-layer}
Let $U=(-1,1)$, $a_\varepsilon=\sqrt\varepsilon$, and
\[
 u_\varepsilon(x)=q\left(\frac{a_\varepsilon-|x|}\varepsilon\right),
 \qquad q(t)=\tanh(t/\sqrt2).
\]
These $H^1$ functions converge to $-1$ in $L^1(U)$, but their energy measures converge to $2c_W\delta_0$. Indeed, on each half interval the change of variables $t=(a_\varepsilon-|x|)/\varepsilon$ gives energy tending to $\int_{\mathbb R}(q')^2=c_W$. Outside any neighborhood of zero the energy tends to zero by exponential decay. The positive-phase interval has length $2a_\varepsilon\to0$, and pointwise convergence holds away from zero, proving the phase convergence by bounded convergence. The limiting phase has perimeter zero. These are competitors, not claimed solutions of the Euler equation; the example explains why a phase limit alone does not record all interface multiplicity.
\end{example}
